\documentclass[11pt]{article}
\usepackage[letterpaper,margin=1in]{geometry}
\usepackage{amsmath,amssymb,amsthm}
\makeatletter
\def\ps@proofguide{%
  \def\@oddhead{\small\sffamily POLYNOMIAL CONTINUED FRACTIONS\hfil Catalan's constant}%
  \def\@oddfoot{\hfil\small\thepage}%
}
\makeatother
\pagestyle{proofguide}
\setlength{\headheight}{14pt}
\setlength{\parindent}{0pt}
\setlength{\parskip}{5pt}

\newtheorem*{theorem}{Identity}
\title{A continued fraction for \(1/(2G)\)}
\author{}
\date{}

\begin{document}
\maketitle
\thispagestyle{proofguide}
\vspace{-3em}

\begin{theorem}
Let \(G=\sum_{k=0}^{\infty}(-1)^k/(2k+1)^2\) be Catalan's constant.
For \(a_n=3n^2+3n+1\) and \(b_n=-2n^4\), the continued fraction satisfies
\[
  a_0+\cfrac{b_1}{a_1+\cfrac{b_2}{a_2+\ddots}}=\frac{1}{2G}.
\]
\end{theorem}

\section{Constructing the moments}

The proof represents a linear combination of the numerator and denominator
continuants as a bounded moment. After normalization, the numerator grows
without bound, so the coefficient of the numerator determines the limit.

For \(k,n\ge0\), put
\[
 d_k=\frac{\binom{2k}{k}}{4^k},\qquad
 c_k=d_k\int_0^{\pi/2}\sin^{2k}x\,dx,\qquad
 B_{n,k}=\int_0^1(1-t)^n t^{2k}\,dt,\qquad
 M_n=\sum_{k\ge0} c_k B_{n,k}.
\]
Integration by parts gives the even sine integral \((\pi/2)d_k\).
Thus \(c_k=(\pi/2)d_k^2\), and the ratio of successive central binomial
coefficients gives
\begin{equation}
 (2k+2)^2 c_{k+1}=(2k+1)^2 c_k.
 \label{eq:coeff}
\end{equation}
The Wallis product \(W_k=\prod_{j=1}^k(2j)^2/((2j-1)(2j+1))\)
satisfies \((2k+1)c_kW_k=\pi/2\), by cancellation in the product.
As \(W_k\to\pi/2\), both \((2k+1)c_k\) and \(2kc_k\) tend to \(1\).

Integration by parts, or the factorial expression
\(B_{n,k}=n!(2k)!/(2k+n+1)!\), gives
\begin{equation}
 B_{0,k}=\frac{1}{2k+1},\qquad
 B_{n+1,k}=\frac{n+1}{2k+n+2}B_{n,k}.
 \label{eq:beta}
\end{equation}
Since \(0\le B_{n,k}\le B_{0,k}\) and \(c_k/(2k+1)=O(k^{-2})\),
all the moment sums converge and \(0\le M_n\le M_0\).

\newpage
\section{The first two moments}

The binomial series for \((1-s^2)^{-1/2}\), integrated from \(0\) to \(s\)
and divided by \(s\), gives
\[
 \frac{\arcsin s}{s}=\sum_{k\ge0}\frac{d_k s^{2k}}{2k+1},
 \qquad 0<s<1.
\]
Set \(s=\sin x\) and integrate on \([0,\pi/2]\). The series terms are
nonnegative, so monotone convergence allows termwise integration.
The resulting sum is \(M_0\). The substitution \(x=2\arctan t\) gives
\[
 M_0=\int_0^{\pi/2}\frac{x}{\sin x}\,dx
     =2\int_0^1\frac{\arctan t}{t}\,dt
     =2\sum_{k\ge0}\frac{(-1)^k}{(2k+1)^2}=2G.
\]
For the last equality, expand \(\arctan t/t\) on \(0<t<1\). The sum of
the integrals of the absolute values is \(\sum_{k\ge0}(2k+1)^{-2}\),
which justifies the exchange of sum and integral.

Equation~\eqref{eq:beta} gives \(B_{0,k}-B_{1,k}=1/(2k+2)\).
Equation~\eqref{eq:coeff} rewrites \(c_k/(2k+2)\) as
\(2(k+1)c_{k+1}-2kc_k\). The sum telescopes, yielding
\[
 M_0-M_1=\lim_{N\to\infty}2Nc_N=1.
\]

\section{The moment recurrence}

Fix \(n\ge0\), and write \(A=(n+2)^2\), \(D=a_{n+1}\), and
\(C=2(n+1)^2\). Define \(f_k=2kc_k(n+1)B_{n,k}\). Using
equations~\eqref{eq:coeff} and~\eqref{eq:beta} together with
\[
 \frac{B_{n,k+1}}{B_{n,k}}
 =\frac{(2k+1)(2k+2)}{(2k+n+2)(2k+n+3)},
\]
substitution and clearing denominators gives
\[
 c_k\bigl(A B_{n+2,k}-D B_{n+1,k}+C B_{n,k}\bigr)=f_k-f_{k+1}.
\]
The boundary terms vanish as \(f_0=0\) and \(f_k\to0\), since
\(2kc_k\to1\) and \(B_{n,k}\le1/(2k+1)\). Summing establishes
\begin{equation}
 (n+2)^2 M_{n+2}=a_{n+1}M_{n+1}-2(n+1)^2M_n .
 \label{eq:recurrence}
\end{equation}

\newpage
\section{Matching the continuants and taking the limit}

For \(m\ge1\), the numerator \(P_m\) and denominator \(Q_m\) of the
\(m\)-th convergent satisfy \(Z_m=a_m Z_{m-1}+b_m Z_{m-2}\).
Begin with \((P_{-1},P_0)=(1,1)\) and
\((Q_{-1},Q_0)=(0,1)\).
Set \(U_n=P_{n-1}/(n!)^2\) and \(V_n=Q_{n-1}/(n!)^2\).
Dividing the recurrence at \(m=n+1\) by \(((n+1)!)^2\) gives
equation~\eqref{eq:recurrence} for both normalized sequences.
Their initial pairs are \((U_0,U_1)=(1,1)\) and \((V_0,V_1)=(0,1)\).
The first two moments therefore determine the entire sequence as
\[
 M_n=M_0U_n-V_n.
\]

To show that \(U_n\) grows, rewrite its recurrence as
\[
 (n+2)^2(U_{n+2}-U_{n+1})
 =2(n+1)^2(U_{n+1}-U_n)+(n+1)U_{n+1}.
\]
The first increment from \(U_1\) to \(U_2=5/4\) is \(1/4\). If
\(U_{n+1}-U_n\ge1/4\) and \(U_{n+1}\ge1\) for \(n\ge1\), the
right-hand side is at least \((n+1)^2/2+(n+1)\ge(n+2)^2/4\).
Induction gives \(U_n\ge1+(n-1)/4\) for \(n\ge1\).

Since \(M_n\) is bounded and \(U_n\to\infty\),
\[
 \frac{V_n}{U_n}=M_0-\frac{M_n}{U_n}\longrightarrow M_0=2G.
\]
The limit is positive because \(M_0-M_1=1\) and \(M_1\ge0\).
Consequently \(V_n\) is eventually nonzero. Inverting the ratio cancels
the factorial scales and gives \(P_{n-1}/Q_{n-1}\to1/(2G)\).

\section*{Reading the Lean proof}

\begin{raggedright}
In \texttt{PcfProof/Catalan.lean}, \texttt{ellCoeff} is \(c_k\) and
\texttt{betaEntry} defines \(B_{n,k}\) by its recurrence.
\texttt{integral\_eq\_catalan} evaluates the integral for \(M_0\),
\texttt{moment\_recurrence} proves equation~\eqref{eq:recurrence}, and
\texttt{limit\_of\_moments} derives convergence.

The theorem
\texttt{CatalanPCF.catalan\_pcf\_target} assembles these steps.
\end{raggedright}
\end{document}
